\documentclass[10pt]{beamer}

\usepackage[utf8]{inputenc}
\usepackage[T1]{fontenc}
\usepackage[spanish]{babel}
\usepackage{amsmath,amssymb,amsthm}
\usepackage{graphicx}
\usepackage{hyperref}

\usetheme{Madrid}
\usecolortheme{default}
\setbeamertemplate{navigation symbols}{}
\setbeamertemplate{itemize items}[circle]
\setbeamersize{text margin left=1em, text margin right=1em}
\makeatletter
\setlength{\abovedisplayskip}{5pt}
\setlength{\belowdisplayskip}{5pt}
\setlength{\abovedisplayshortskip}{3pt}
\setlength{\belowdisplayshortskip}{3pt}
\makeatother

\setbeamerfont{title}{size=\large}
\setbeamerfont{frametitle}{size=\normalsize}
\setbeamerfont{author}{size=\small}
\setbeamerfont{institute}{size=\tiny}
\setbeamerfont{date}{size=\tiny}
\setbeamerfont{normal text}{size=\footnotesize}

\title[Igualdad de Parseval]{\fontsize{18}{20}\selectfont Igualdad de Parseval}
\author[\fontsize{5}{6}\selectfont Lester Hernández \and Emilio Reyes]{\fontsize{11}{13}\selectfont Lester Iván Hernández Hernández - 20242000637 \\ Emilio Gabriel Reyes Pérez - 20242000121}
\institute[MM-415]{\fontsize{12}{14}\selectfont Ecuaciones Diferenciales Parciales (MM-415)}
\date[\today]{\fontsize{10}{12}\selectfont \today}

\begin{document}

\begin{frame}[plain]
  \maketitle
\end{frame}

\begin{frame}{Enunciado del problema}
Sea $f(x) \in L^2[-p,p] \longrightarrow \mathbb{R}$ cuyo desarrollo de Fourier está dado por
\[
  f(x) = \frac{a_0}{2} + \sum_{n=1}^{\infty} \left(a_n \cos\frac{n\pi}{p}x + b_n \sin\frac{n\pi}{p}x\right).
\]

Probar la igualdad de Parseval:
\[
  \langle f,f\rangle = \|f\|^2 = \int_{-p}^{p} f(x)^2\,dx
  = p\left(\frac{a_0^2}{2} + \sum_{n=1}^{\infty}(a_n^2+b_n^2)\right).
\]
\end{frame}

\begin{frame}{Definición de convergencia en $L^2$}
Una sucesión de funciones $\{f_n\}$ converge a una función $f$ en la norma $L^2$ si
\[
  \lim_{n\to\infty} \int_{-p}^{p} |f_n(x)-f(x)|^2\,dx = 0.
\]

Esto significa que la energía o la norma cuadrática de la diferencia tiende a cero.

En este caso, la norma de $f$ se define como
\[
  \|f\|^2 = \int_{-p}^{p} f(x)^2\,dx.
\]
\end{frame}

\begin{frame}{Solución}
Siguiendo la definición de norma con función peso $1$,
\[
  \|f\|^2 = \int_{-p}^{p} f(x)^2\,dx
  = \int_{-p}^{p} \left(\frac{a_0}{2} + \sum_{n=1}^{\infty}\left(a_n\cos\frac{n\pi}{p}x + b_n\sin\frac{n\pi}{p}x\right)\right)^2 dx.
\]

Ahora expandimos el cuadrado.
\end{frame}

\begin{frame}{Paso 1: desarrollo del cuadrado}
\[
\small
\begin{aligned}
\int_{-p}^{p} f(x)^2\,dx
= \int_{-p}^{p} \Bigg[ &\frac{a_0^2}{4}
+ a_0\sum_{n=1}^{\infty}\left(a_n\cos\frac{n\pi}{p}x+b_n\sin\frac{n\pi}{p}x\right) \\
&+ \left(\sum_{n=1}^{\infty}\left(a_n\cos\frac{n\pi}{p}x+b_n\sin\frac{n\pi}{p}x\right)\right)^2 \Bigg] dx.
\end{aligned}
\]

Aplicando la linealidad de la integral:
\[
\small
\int_{-p}^{p} \frac{a_0^2}{4}\,dx
+ \int_{-p}^{p} a_0\sum_{n=1}^{\infty}\left(\cdots\right)dx
+ \int_{-p}^{p}\left(\sum_{n=1}^{\infty}\left(\cdots\right)\right)^2 dx.
\]
\end{frame}

\begin{frame}{Paso 2: primer término}
\[
\int_{-p}^{p} \frac{a_0^2}{4}\,dx
= \frac{a_0^2}{4}\int_{-p}^{p}dx
= \frac{a_0^2}{4}(2p)
= \frac{a_0^2}{2}p.
\]

Además,
\[
\int_{-p}^{p} a_0\sum_{n=1}^{\infty}\left(a_n\cos\frac{n\pi}{p}x+b_n\sin\frac{n\pi}{p}x\right)dx = 0,
\]
porque
\[
\int_{-p}^{p} \cos\frac{n\pi}{p}x\,dx = 0,
\qquad
\int_{-p}^{p} \sin\frac{n\pi}{p}x\,dx = 0.
\]
\end{frame}

\begin{frame}{Paso 3: ortogonalidad de senos y cosenos}
Consideremos ahora el término cuadrático:
\[
\begin{aligned}
&\int_{-p}^{p}\left(\sum_{n=1}^{\infty} a_n\cos\frac{n\pi}{p}x + b_n\sin\frac{n\pi}{p}x\right)^2dx \\
&= \int_{-p}^{p}\sum_{n=1}^{\infty}\sum_{m=1}^{\infty}\bigg(
  a_n a_m \cos\frac{n\pi}{p}x\cos\frac{m\pi}{p}x
  + a_n b_m \cos\frac{n\pi}{p}x\sin\frac{m\pi}{p}x \\
&\hspace{3.2cm}+ b_n a_m \sin\frac{n\pi}{p}x\cos\frac{m\pi}{p}x
  + b_n b_m \sin\frac{n\pi}{p}x\sin\frac{m\pi}{p}x\bigg)dx.
\end{aligned}
\]

Por ortogonalidad,
\[
\int_{-p}^{p}\cos\frac{n\pi}{p}x\cos\frac{m\pi}{p}x\,dx = 0 \quad (n\neq m),
\]
\[
\int_{-p}^{p}\sin\frac{n\pi}{p}x\sin\frac{m\pi}{p}x\,dx = 0 \quad (n\neq m),
\]
\[
\int_{-p}^{p}\cos\frac{n\pi}{p}x\sin\frac{m\pi}{p}x\,dx = 0,
\qquad
\int_{-p}^{p}\sin\frac{n\pi}{p}x\cos\frac{m\pi}{p}x\,dx = 0.
\]
\end{frame}

\begin{frame}{Paso 4: simplificación final}
Entonces solo sobreviven los términos con $n=m$:
\[
\begin{aligned}
&\int_{-p}^{p}\left(\sum_{n=1}^{\infty} a_n\cos\frac{n\pi}{p}x + b_n\sin\frac{n\pi}{p}x\right)^2dx \\
&= \sum_{n=1}^{\infty}\left(a_n^2\int_{-p}^{p}\cos^2\frac{n\pi}{p}x\,dx + b_n^2\int_{-p}^{p}\sin^2\frac{n\pi}{p}x\,dx\right).
\end{aligned}
\]

Usando que
\[
\int_{-p}^{p}\cos^2\frac{n\pi}{p}x\,dx = p,
\qquad
\int_{-p}^{p}\sin^2\frac{n\pi}{p}x\,dx = p,
\]
obtenemos
\[
\sum_{n=1}^{\infty}\left(a_n^2p+b_n^2p\right)=p\sum_{n=1}^{\infty}(a_n^2+b_n^2).
\]
\end{frame}

\begin{frame}{Conclusión}
Reuniendo todos los términos:
\[
\begin{aligned}
\|f\|^2 &= \int_{-p}^{p} f(x)^2dx \\
&= \frac{a_0^2}{2}p + 0 + p\sum_{n=1}^{\infty}(a_n^2+b_n^2) \\
&= p\left(\frac{a_0^2}{2} + \sum_{n=1}^{\infty}(a_n^2+b_n^2)\right).
\end{aligned}
\]

Por lo tanto,
\[
\boxed{\langle f,f\rangle = \|f\|^2 = p\left(\frac{a_0^2}{2} + \sum_{n=1}^{\infty}(a_n^2+b_n^2)\right)}
\]
que es exactamente la igualdad de Parseval.
\end{frame}

\end{document}
